\documentclass[11pt,a4paper]{article}
\usepackage{fontspec}
\usepackage{babel}
\usepackage[margin=2.5cm]{geometry}
\usepackage{enumitem}
\usepackage{booktabs}
\usepackage{xcolor}
\usepackage{hyperref}
\usepackage{pgffor}
\usepackage{tikz}
\usetikzlibrary{calc}

\definecolor{coral}{HTML}{cc785c}
\definecolor{ink}{HTML}{141413}
\definecolor{muted}{HTML}{6c6a64}
\definecolor{done}{HTML}{4caf50}
\definecolor{wip}{HTML}{ff9800}
\definecolor{todo}{HTML}{9e9e9e}

\hypersetup{colorlinks=true, linkcolor=coral, urlcolor=coral}

\babelprovide[import, onchar=ids fonts]{thai}
\babelprovide[import, onchar=ids fonts]{english}

\babelfont{rm}{Noto Sans}
\babelfont[thai]{rm}{Noto Sans Thai}
\babelfont{tt}{Noto Sans Mono}
\babelfont[thai]{tt}{Noto Sans Thai}

\directlua{
  local glyph_id = node.id("glyph")
  local penalty_id = node.id("penalty")

  local function is_thai(c)
    return c >= 0x0E01 and c <= 0x0E5B
  end

  local function is_thai_combining(c)
    return (c >= 0x0E31 and c <= 0x0E3A) or (c >= 0x0E47 and c <= 0x0E4E)
  end

  luatexbase.add_to_callback("pre_linebreak_filter", function(head)
    for n in node.traverse_id(glyph_id, head) do
      local prev = n.prev
      if prev and prev.id == glyph_id and is_thai(prev.char)
         and is_thai(n.char) and not is_thai_combining(n.char) then
        local p = node.new(penalty_id)
        p.penalty = 0
        node.insert_before(head, n, p)
      end
    end
    return head
  end, "thai_break")
}

\emergencystretch=1em

%% ─── Progress bar macro ───────────────────────────────────────
\newcommand{\progressbar}[1]{%
  \begin{tikzpicture}[baseline=-0.5ex]
    \fill[muted!20] (0,0) rectangle (4,0.35);
    \fill[coral]    (0,0) rectangle ({4*#1/100},0.35);
    \node[right] at (4.1,0.175) {\small\color{muted}#1\%};
  \end{tikzpicture}%
}

%% ─── Status badge macros ──────────────────────────────────────
\newcommand{\done}{\textcolor{done}{\textbf{[เสร็จ]}}}
\newcommand{\wip}{\textcolor{wip}{\textbf{[กำลังทำ]}}}
\newcommand{\todo}{\textcolor{todo}{[รอ]}}

\begin{document}

\begin{center}
  {\LARGE \textbf{รายงานความคืบหน้า}}\\[6pt]
  {\large ชื่อโครงการ / Thesis Title}\\[4pt]
  {\color{muted}\small Progress Report · \today}
\end{center}

\noindent\rule{\textwidth}{1.5pt}
\vspace{8pt}

%% ─── Overview ─────────────────────────────────────────────────
\section*{ภาพรวม}

\noindent
\begin{tabular}{@{}ll@{}}
  \textbf{นักศึกษา}    & ชื่อ-นามสกุล \\[3pt]
  \textbf{อาจารย์ที่ปรึกษา} & ชื่ออาจารย์ \\[3pt]
  \textbf{สถาบัน}      & มหาวิทยาลัย / ภาควิชา \\[3pt]
  \textbf{ช่วงเวลา}    & เดือน ปี -- เดือน ปี \\[3pt]
\end{tabular}

\vspace{10pt}
\noindent\textbf{ความคืบหน้าโดยรวม:} \progressbar{35}

\vspace{12pt}
\noindent\rule{\textwidth}{0.4pt}

%% ─── Work done ────────────────────────────────────────────────
\section*{งานที่ทำเสร็จแล้ว (ช่วงนี้)}

\begin{itemize}[nosep, leftmargin=2em]
  \item[\done] งานที่ 1: อธิบายสั้น ๆ
  \item[\done] งานที่ 2: อธิบายสั้น ๆ
  \item[\done] งานที่ 3: อธิบายสั้น ๆ
\end{itemize}

%% ─── In progress ──────────────────────────────────────────────
\section*{งานที่กำลังดำเนินการ}

\begin{itemize}[nosep, leftmargin=2em]
  \item[\wip] งานที่ A: อธิบาย — \progressbar{60}
  \item[\wip] งานที่ B: อธิบาย — \progressbar{30}
\end{itemize}

%% ─── Upcoming ─────────────────────────────────────────────────
\section*{แผนงานต่อไป}

\begin{itemize}[nosep, leftmargin=2em]
  \item[\todo] งานที่ X: กำหนด [วันที่]
  \item[\todo] งานที่ Y: กำหนด [วันที่]
  \item[\todo] งานที่ Z: กำหนด [วันที่]
\end{itemize}

%% ─── Milestone table ──────────────────────────────────────────
\section*{ตาราง Milestone}

\vspace{6pt}
\noindent
\begin{tabular}{@{}lllc@{}}
  \toprule
  \textbf{Milestone} & \textbf{กำหนดส่ง} & \textbf{สถานะ} & \textbf{\%} \\
  \midrule
  Chapter 1 (Introduction)      & มี.ค. 2025  & \done  & 100 \\
  Chapter 2 (Literature Review) & พ.ค. 2025  & \wip   & 60  \\
  Chapter 3 (Methodology)       & ก.ค. 2025  & \todo  & 0   \\
  Chapter 4 (Experiments)       & ก.ย. 2025  & \todo  & 0   \\
  Chapter 5 (Conclusion)        & ต.ค. 2025  & \todo  & 0   \\
  \bottomrule
\end{tabular}

%% ─── Blockers ─────────────────────────────────────────────────
\section*{ปัญหาและอุปสรรค}

\begin{itemize}[nosep, leftmargin=2em]
  \item อุปสรรคที่ 1 — แนวทางแก้ไข
  \item อุปสรรคที่ 2 — แนวทางแก้ไข
\end{itemize}

%% ─── Questions for advisor ────────────────────────────────────
\section*{คำถามสำหรับอาจารย์ที่ปรึกษา}

\begin{enumerate}[nosep, leftmargin=2em]
  \item คำถามที่ 1
  \item คำถามที่ 2
\end{enumerate}

\vspace{16pt}
\noindent\rule{\textwidth}{1.5pt}
\noindent{\color{muted}\small สรุปโดย Claude Sonnet 4.6 · \today}

\end{document}
